% =============================================================================
% Aula 02 — HTML/JS estático + API JSON · Projeto StockSales
% Beamer + tema Madrid | Curso Entra21 — Spring Framework
% Compilar: pdflatex slides-aula-02.tex  (duas vezes se usar TOC)
% =============================================================================
\documentclass[aspectratio=169,11pt]{beamer}

\usetheme{Madrid}
\usecolortheme{default}

\usepackage[T1]{fontenc}
\usepackage[utf8]{inputenc}
\usepackage[brazil]{babel}
\usepackage{booktabs}
\usepackage{graphicx}
\usepackage{hyperref}
\usepackage{listings}
\usepackage{xcolor}
\usepackage{tikz}
\usetikzlibrary{arrows.meta, positioning, shapes.geometric, fit}

% --- Cores e código ---
\definecolor{codebg}{RGB}{245,245,245}
\definecolor{codegreen}{RGB}{40,120,40}
\definecolor{codeblue}{RGB}{30,70,160}
\definecolor{codered}{RGB}{180,50,50}

\lstset{
  basicstyle=\ttfamily\small,
  backgroundcolor=\color{codebg},
  frame=single,
  rulecolor=\color{black!20},
  breaklines=true,
  columns=fullflexible,
  keepspaces=true,
  showstringspaces=false,
  keywordstyle=\color{codeblue}\bfseries,
  commentstyle=\color{codegreen},
  stringstyle=\color{codered},
  tabsize=2,
  literate={á}{{\'a}}1 {é}{{\'e}}1 {í}{{\'i}}1 {ó}{{\'o}}1 {ú}{{\'u}}1
           {Á}{{\'A}}1 {É}{{\'E}}1 {Í}{{\'I}}1 {Ó}{{\'O}}1 {Ú}{{\'U}}1
           {ã}{{\~a}}1 {õ}{{\~o}}1 {Ã}{{\~A}}1 {Õ}{{\~O}}1
           {ç}{{\c{c}}}1 {Ç}{{\c{C}}}1
           {←}{{$\leftarrow$}}1 {→}{{$\rightarrow$}}1
}

\title{Aula 02 --- Spring Framework}
\subtitle{HTML/JS + API JSON · Projeto StockSales}
\author{Mauricio Duailibi Neto}
\institute{Turma Java Entra21}
\date{}

\begin{document}

% =============================================================================
\begin{frame}
  \titlepage
\end{frame}

% =============================================================================
\begin{frame}{Como será a aula de hoje}
  Duas partes, com intervalo no meio:

  \vspace{0.8em}
  \begin{block}{Parte 1 --- 18h30 às 20h}
    \begin{itemize}
      \item Revisão da Aula 01
      \item Estrutura de um projeto Spring Boot
      \item Como o HTML conversa com o backend
      \item Exercícios guiados
    \end{itemize}
  \end{block}

  \vspace{0.4em}
  \begin{block}{Parte 2 --- 20h20 às 22h}
    \begin{itemize}
      \item Apresentação do projeto do curso: \textbf{StockSales}
      \item Configuração inicial + primeira página + listar produtos
    \end{itemize}
  \end{block}
\end{frame}

% =============================================================================
\begin{frame}{Mapa da aula}
  \tableofcontents
\end{frame}

% #############################################################################
\section{Revisão da Aula 01}

% =============================================================================
\begin{frame}{O que construímos na Aula 01}
  Na última aula, o foco foi fazer o Spring responder no navegador:

  \vspace{0.6em}
  \begin{enumerate}
    \item Instalar JDK 21 + VS Code + extensões
    \item Gerar um projeto no \texttt{start.spring.io}
    \item Criar um \texttt{@RestController}
    \item Responder texto em endereços como \texttt{/} e \texttt{/ola}
    \item Usar \texttt{@RequestParam}, por exemplo \texttt{/saudacao?nome=Ana}
  \end{enumerate}

  \vspace{0.6em}
  \begin{alertblock}{Lembrete}
    O navegador pede uma URL $\rightarrow$ o Controller responde.
  \end{alertblock}
\end{frame}

% =============================================================================
\begin{frame}{Anotações que precisam ficar na cabeça}
  \begin{description}
    \item[\texttt{@RestController}] Diz: ``esta classe recebe pedidos HTTP''.
    \item[\texttt{@GetMapping("/caminho")}] Diz: ``quando alguém acessar
      este caminho com GET, rode este método''.
    \item[\texttt{@RequestParam}] Lê valores da URL depois do \texttt{?},
      por exemplo \texttt{?nome=Ana}.
    \item[\texttt{pom.xml}] Lista de dependências do Maven,
      por exemplo Spring Web.
  \end{description}

  \vspace{0.6em}
  Sem isso, o restante da disciplina fica mais difícil.\\
  Com isso, vocês já conseguem conversar com o servidor.
\end{frame}

% =============================================================================
\begin{frame}{Fluxo que já vimos}
  \begin{center}
  \begin{tikzpicture}[
      box/.style={draw, rounded corners, align=center, minimum width=2.5cm,
                  minimum height=1.0cm, font=\small\bfseries},
      seta/.style={-{Latex}, thick}
    ]
    \node[box, fill=blue!15] (nav) {Navegador};
    \node[box, fill=orange!20, right=1.0cm of nav] (boot) {Spring Boot};
    \node[box, fill=green!15, right=1.0cm of boot] (ctrl) {Controller};
    \node[box, fill=yellow!25, below=0.85cm of ctrl] (txt) {Texto\\{\normalfont\tiny resposta}};

    \draw[seta] (nav) -- node[above, font=\tiny]{GET /ola} (boot);
    \draw[seta] (boot) -- (ctrl);
    \draw[seta] (ctrl) -- (txt);
    \draw[seta] (txt.west) -| ++(-1.0,0) |- (nav.south);
  \end{tikzpicture}
  \end{center}

  \vspace{0.8em}
  Isso funciona --- mas a resposta era \textbf{texto puro}.\\
  Hoje vamos dar o próximo passo: \textbf{páginas de verdade}.
\end{frame}

% =============================================================================
\begin{frame}{O limite do texto puro}
  Na Aula 01, o navegador mostrava algo assim:

  \vspace{0.5em}
  \begin{center}
    \texttt{Olá, Spring Boot! Bem-vindo à disciplina.}
  \end{center}

  \vspace{0.8em}
  Isso é ótimo para aprender o fluxo.\\
  Mas um sistema real precisa de:

  \vspace{0.4em}
  \begin{itemize}
    \item Layout: cores, títulos, botões
    \item Listas, tabelas, formulários
    \item Interação: clicar e atualizar a tela
  \end{itemize}

  \vspace{0.6em}
  Para isso, usamos \textbf{HTML + CSS + JavaScript}.
\end{frame}

% =============================================================================
\begin{frame}{Pergunta-guia da Aula 02}
  \begin{center}
    \Large
    Como o navegador mostra uma \textbf{página}\\
    e, ao mesmo tempo, busca \textbf{dados}\\
    que vêm do Spring?
  \end{center}

  \vspace{1.2em}
  A resposta curta:\\
  \textbf{HTML/JS estático} + \textbf{API JSON}.

  \vspace{0.8em}
  Vamos entender isso com calma.
\end{frame}

% #############################################################################
\section{Arquitetura de um projeto Spring Boot}

% =============================================================================
\begin{frame}{Antes de HTML: onde cada coisa mora}
  Um projeto Spring Boot não é ``um arquivo só''.\\
  É uma \textbf{casinha organizada}.

  \vspace{0.8em}
  Se vocês souberem \textbf{onde procurar},\\
  programar fica muito menos assustador.

  \vspace{0.8em}
  \begin{block}{Meta deste bloco}
    Olhar a estrutura de pastas e saber dizer:\\
    ``isso é código Java'' · ``isso é página'' · ``isso é dependência''.
  \end{block}
\end{frame}

% =============================================================================
\begin{frame}[fragile]{Estrutura típica do projeto}
\begin{lstlisting}[basicstyle=\ttfamily\footnotesize]
meu-projeto/
├── pom.xml
├── src/
│   └── main/
│       ├── java/
│       │   └── com/entra21/...
│       │       ├── App.java
│       │       └── ...Controller.java
│       └── resources/
│           ├── application.properties
│           └── static/          ← HTML, CSS, JS
└── target/                      ← gerado ao compilar
\end{lstlisting}

  \vspace{0.3em}
  Hoje a pasta nova mais importante é: \texttt{static/}.
\end{frame}

% =============================================================================
\begin{frame}{O que cada pasta faz}
  \begin{description}
    \item[\texttt{pom.xml}] ``Lista de compras'' do projeto:
      quais bibliotecas baixar.
    \item[\texttt{src/main/java}] Código Java: Controllers, regras, etc.
    \item[\texttt{src/main/resources}] Configuração e arquivos da aplicação.
    \item[\texttt{static/}] Páginas e estilos que o navegador baixa:
      \texttt{.html}, \texttt{.css}, \texttt{.js}, imagens.
    \item[\texttt{target/}] Resultado da compilação --- \textbf{não editamos}
      na mão.
  \end{description}
\end{frame}

% =============================================================================
\begin{frame}{Duas portas no mesmo servidor}
  No Spring Boot, o mesmo programa --- porta 8080 --- pode servir:

  \vspace{0.6em}
  \begin{columns}[T]
    \begin{column}{0.48\textwidth}
      \begin{block}{Páginas}
        \texttt{/}\\
        \texttt{/index.html}\\
        \texttt{/css/estilo.css}\\
        \texttt{/js/app.js}
      \end{block}
      Arquivos da pasta \texttt{static/}.
    \end{column}
    \begin{column}{0.48\textwidth}
      \begin{block}{API --- dados}
        \texttt{/api/produtos}\\
        \texttt{/api/clientes}\\
        \texttt{/api/vendas}
      \end{block}
      Métodos do \texttt{@RestController}.
    \end{column}
  \end{columns}

  \vspace{0.8em}
  \begin{alertblock}{Frase de ouro}
    A página é a \textbf{vitrine}.\\
    A API é o \textbf{estoque de dados} por trás da vitrine.
  \end{alertblock}
\end{frame}

% =============================================================================
\begin{frame}{Analogia do restaurante}
  \begin{itemize}
    \item \textbf{HTML} = o cardápio e a mesa: o que o cliente vê
    \item \textbf{CSS} = a decoração: cores, fontes, espaçamento
    \item \textbf{JavaScript} = o garçom: vai à cozinha e traz o pedido
    \item \textbf{API Spring} = a cozinha: prepara e devolve os dados
    \item \textbf{JSON} = a bandeja com o prato organizado
  \end{itemize}

  \vspace{0.8em}
  O usuário quase nunca ``entra na cozinha''.\\
  Ele conversa com a interface --- e a interface fala com a API.
\end{frame}

% =============================================================================
\begin{frame}{Diagrama: página e API juntas}
  \begin{center}
  \begin{tikzpicture}[
      box/.style={draw, rounded corners, align=center, minimum width=2.3cm,
                  minimum height=0.95cm, font=\small\bfseries},
      seta/.style={-{Latex}, thick}
    ]
    \node[box, fill=blue!15] (nav) {Navegador};
    \node[box, fill=orange!20, right=1.3cm of nav] (boot) {Spring\\Boot};
    \node[box, fill=green!20, above right=0.6cm and 1.0cm of boot] (st) {static/\\HTML CSS JS};
    \node[box, fill=yellow!30, below right=0.6cm and 1.0cm of boot] (api) {Controller\\/api/...};

    \draw[seta] (nav) -- node[above, font=\tiny]{1) GET /} (boot);
    \draw[seta] (boot) -- (st);
    \draw[seta] (nav.east) to[bend left=12] node[above, font=\tiny]{2) GET /api/...} (api.west);
    \draw[seta] (api) -- node[below, font=\tiny]{JSON} (boot);
  \end{tikzpicture}
  \end{center}

  \vspace{0.9em}
  \begin{enumerate}
    \item O navegador baixa a página \texttt{index.html}.
    \item O JavaScript da página chama a API e recebe dados.
  \end{enumerate}
\end{frame}

% =============================================================================
\begin{frame}{Por que separar página e API}
  Porque isso facilita a vida no dia a dia:

  \vspace{0.5em}
  \begin{itemize}
    \item A tela cuida de mostrar as informações
    \item O backend cuida de devolver os dados
    \item Cada parte fica mais fácil de entender e testar
    \item O mesmo backend pode alimentar várias telas depois
  \end{itemize}

  \vspace{0.7em}
  \begin{block}{No nosso curso}
    Vamos guardar o HTML/CSS/JS na pasta \texttt{static/}\\
    e os dados nos Controllers em caminhos \texttt{/api/...}.
  \end{block}
\end{frame}

% =============================================================================
\begin{frame}{Classes Java que vão aparecer bastante}
  \begin{description}
    \item[Application] Classe com \texttt{main} --- sobe o Spring.
    \item[Controller] Recebe o pedido HTTP e devolve a resposta.
    \item[Classe de dados] Representa algo do sistema,
      por exemplo \texttt{Produto} ou \texttt{Aluno}.
  \end{description}

  \vspace{0.8em}
  Hoje: \textbf{Application + Controller + classe de dados + HTML/JS}.\\
  Sem banco de dados ainda.
\end{frame}

% =============================================================================
\begin{frame}{E o banco de dados}
  \begin{center}
    \Large Ainda \textbf{não}.
  \end{center}

  \vspace{0.8em}
  Nas primeiras aulas do StockSales, os dados ficam\\
  em uma \textbf{lista na memória} do Java.

  \vspace{0.6em}
  \begin{itemize}
    \item Reiniciou o servidor? Os dados voltam ao que estava no código.
    \item Isso é proposital: aprendemos a API e a tela primeiro.
    \item Banco entra mais à frente no curso.
  \end{itemize}
\end{frame}

% #############################################################################
\section{HTML estático + API JSON}

% =============================================================================
\begin{frame}{O que é HTML estático no Spring}
  ``Estático'' aqui significa:

  \vspace{0.5em}
  \begin{itemize}
    \item O arquivo já existe em \texttt{src/main/resources/static/}
    \item O Spring \textbf{entrega o arquivo} para o navegador
    \item Não precisamos criar um Controller só para servir o HTML
  \end{itemize}

  \vspace{0.8em}
  Exemplos:

  \begin{center}
  \begin{tabular}{ll}
    \toprule
    Arquivo & URL \\
    \midrule
    \texttt{static/index.html} & \texttt{http://localhost:8080/} \\
    \texttt{static/produtos.html} & \texttt{.../produtos.html} \\
    \texttt{static/css/estilo.css} & \texttt{.../css/estilo.css} \\
    \texttt{static/js/app.js} & \texttt{.../js/app.js} \\
    \bottomrule
  \end{tabular}
  \end{center}
\end{frame}

% =============================================================================
\begin{frame}[fragile]{Exemplo mínimo de index.html}
\begin{lstlisting}[basicstyle=\ttfamily\footnotesize]
<!DOCTYPE html>
<html lang="pt-BR">
<head>
  <meta charset="UTF-8">
  <title>StockSales</title>
  <link rel="stylesheet" href="/css/estilo.css">
</head>
<body>
  <h1>StockSales</h1>
  <p>Bem-vindo ao sistema da turma!</p>
  <ul id="lista-produtos"></ul>

  <script src="https://cdn.jsdelivr.net/npm/axios/dist/axios.min.js"></script>
  <script src="/js/app.js"></script>
</body>
</html>
\end{lstlisting}
\end{frame}

% =============================================================================
\begin{frame}{O que é JSON}
  JSON = jeito padrão de \textbf{organizar dados em texto}\\
  para programas trocarem informação.

  \vspace{0.6em}
  Parece um objeto Java, mas é texto:

  \vspace{0.4em}
  \begin{center}
  \texttt{\{"nome": "Teclado", "preco": 150.0, "estoque": 10\}}
  \end{center}

  \vspace{0.8em}
  \begin{itemize}
    \item Chaves entre aspas: \texttt{"nome"}
    \item Dois-pontos separa chave e valor
    \item Texto também entre aspas; números sem aspas
    \item Lista: \texttt{[ {...}, {...} ]}
  \end{itemize}
\end{frame}

% =============================================================================
\begin{frame}[fragile]{JSON de uma lista de produtos}
\begin{lstlisting}[basicstyle=\ttfamily\footnotesize]
[
  {
    "id": 1,
    "nome": "Teclado Mecanico",
    "preco": 250.0,
    "estoque": 12
  },
  {
    "id": 2,
    "nome": "Mouse Gamer",
    "preco": 120.0,
    "estoque": 30
  }
]
\end{lstlisting}

  O JavaScript sabe ler isso.\\
  O Spring gera isso a partir dos objetos Java.
\end{frame}

% =============================================================================
\begin{frame}[fragile]{No Java: a classe Produto}
  Criamos uma classe normal, com atributos e \texttt{get/set}:

\begin{lstlisting}[language=Java,basicstyle=\ttfamily\footnotesize]
public class Produto {
  private Long id;
  private String nome;
  private double preco;
  private int estoque;

  public Produto(Long id, String nome,
                 double preco, int estoque) {
    this.id = id;
    this.nome = nome;
    this.preco = preco;
    this.estoque = estoque;
  }

  public Long getId() { return id; }
  public String getNome() { return nome; }
  public double getPreco() { return preco; }
  public int getEstoque() { return estoque; }
}
\end{lstlisting}
\end{frame}

% =============================================================================
\begin{frame}[fragile]{Por que os métodos get}
  O Spring monta o JSON usando os \texttt{get...}.

  \vspace{0.4em}
  \begin{center}
  \begin{tabular}{ll}
    \toprule
    Método Java & Campo no JSON \\
    \midrule
    \texttt{getNome()} & \texttt{"nome"} \\
    \texttt{getPreco()} & \texttt{"preco"} \\
    \texttt{getEstoque()} & \texttt{"estoque"} \\
    \bottomrule
  \end{tabular}
  \end{center}

  \vspace{0.6em}
  Sem o \texttt{get}, o campo quase não aparece no JSON.\\
  Por isso: atributos \texttt{private} + \texttt{get} públicos.
\end{frame}

% =============================================================================
\begin{frame}[fragile]{Controller que devolve JSON}
\begin{lstlisting}[language=Java,basicstyle=\ttfamily\footnotesize]
@RestController
public class ProdutoController {

  @GetMapping("/api/produtos")
  public List<Produto> listar() {
    List<Produto> produtos = new ArrayList<>();
    produtos.add(new Produto(1L, "Teclado", 250.0, 12));
    produtos.add(new Produto(2L, "Mouse", 120.0, 30));
    return produtos;
  }
}
\end{lstlisting}

  \vspace{0.3em}
  Importes que costumam faltar:\\
  \texttt{import java.util.ArrayList;}\\
  \texttt{import java.util.List;}

  \vspace{0.3em}
  Teste: \texttt{http://localhost:8080/api/produtos}
\end{frame}

% =============================================================================
\begin{frame}{Por que o caminho começa com /api}
  É uma \textbf{convenção} nossa --- um combinado claro:

  \vspace{0.5em}
  \begin{itemize}
    \item \texttt{/api/...} $\rightarrow$ dados em JSON
    \item \texttt{/}, \texttt{/produtos.html} $\rightarrow$ páginas
  \end{itemize}

  \vspace{0.8em}
  Assim ninguém confunde:
  \begin{itemize}
    \item ``Estou pedindo a página?''
    \item ``Ou estou pedindo os dados?''
  \end{itemize}
\end{frame}

% =============================================================================
\begin{frame}{Como o JS fala com a API}
  Precisamos de uma ferramenta no JavaScript que:
  \begin{enumerate}
    \item Chame um endereço \texttt{/api/...}
    \item Receba o JSON
    \item Entregue esses dados para a gente usar na tela
  \end{enumerate}

  \vspace{0.7em}
  No curso vamos usar o \textbf{Axios}.

  \vspace{0.4em}
  \begin{block}{Em português}
    Axios = biblioteca pronta para fazer pedidos HTTP\\
    de um jeito simples e organizado.
  \end{block}
\end{frame}

% =============================================================================
\begin{frame}[fragile]{Como incluir o Axios na página}
  Não vamos instalar nada no computador.\\
  Colocamos o Axios pela internet, com uma tag \texttt{script}:

\begin{lstlisting}[basicstyle=\ttfamily\footnotesize]
<script src="https://cdn.jsdelivr.net/npm/axios/dist/axios.min.js"></script>
<script src="/js/app.js"></script>
\end{lstlisting}

  \vspace{0.5em}
  \begin{alertblock}{Ordem importa}
    1) Axios primeiro\\
    2) Nosso \texttt{app.js} depois\\[0.3em]
    Se inverter, o \texttt{app.js} tenta usar o Axios\\
    antes dele existir --- e quebra.
  \end{alertblock}
\end{frame}

% =============================================================================
\begin{frame}[fragile]{Primeiro pedido com Axios}
  Buscar a lista de produtos:

\begin{lstlisting}[basicstyle=\ttfamily\footnotesize]
axios.get("/api/produtos")
  .then(function (resposta) {
    console.log(resposta.data);
  });
\end{lstlisting}

  \vspace{0.5em}
  \begin{description}
    \item[\texttt{axios.get}] Pede dados com GET.
    \item[\texttt{resposta.data}] Já vem o JSON convertido.\\
      Não precisamos chamar \texttt{.json()} na mão.
  \end{description}

  \vspace{0.4em}
  \begin{block}{Leitura em voz alta}
    ``Axios, busca \texttt{/api/produtos}.\\
    Quando voltar, me mostra \texttt{resposta.data}.''
  \end{block}
\end{frame}

% =============================================================================
\begin{frame}[fragile]{Montando a lista na tela}
\begin{lstlisting}[basicstyle=\ttfamily\footnotesize]
var ul = document.getElementById("lista-produtos");

axios.get("/api/produtos")
  .then(function (resposta) {
    var produtos = resposta.data;
    ul.innerHTML = "";
    for (var i = 0; i < produtos.length; i++) {
      var p = produtos[i];
      var li = document.createElement("li");
      li.textContent = p.nome + " - R$ " + p.preco
        + " - estoque: " + p.estoque;
      ul.appendChild(li);
    }
  });
\end{lstlisting}
\end{frame}

% =============================================================================
\begin{frame}{Lembra do @RequestParam?}
  Na Aula 01, o backend já lia parâmetros na URL:

  \vspace{0.4em}
  \begin{center}
    \texttt{/saudacao?nome=Ana}
  \end{center}

  \vspace{0.6em}
  O \texttt{?} começa a lista de parâmetros.\\
  O \texttt{\&} separa vários parâmetros:

  \vspace{0.4em}
  \begin{center}
    \texttt{/soma?a=2\&b=3}
  \end{center}

  \vspace{0.6em}
  Agora vamos aprender a \textbf{enviar} esses parâmetros\\
  a partir do JavaScript, com Axios.
\end{frame}

% =============================================================================
\begin{frame}{O que é ``enviar parâmetros''}
  Às vezes a API precisa de informação extra, por exemplo:

  \vspace{0.5em}
  \begin{itemize}
    \item Qual o \textbf{nome} da pessoa?
    \item Qual o \textbf{id} do produto?
    \item Qual o texto da \textbf{busca}?
  \end{itemize}

  \vspace{0.6em}
  Esses valores vão na URL, depois do \texttt{?}.

  \vspace{0.5em}
  \begin{block}{Exemplo}
    \texttt{/api/saudacao?nome=Ana}\\
    O parâmetro se chama \texttt{nome}.\\
    O valor é \texttt{Ana}.
  \end{block}
\end{frame}

% =============================================================================
\begin{frame}[fragile]{Backend: receber o parâmetro}
  No Controller, igual à Aula 01:

\begin{lstlisting}[language=Java,basicstyle=\ttfamily\footnotesize]
@GetMapping("/api/saudacao")
public Mensagem saudacao(
    @RequestParam String nome) {
  return new Mensagem("Ola, " + nome + "!");
}
\end{lstlisting}

  \vspace{0.4em}
  Teste manual no navegador:\\
  \texttt{http://localhost:8080/api/saudacao?nome=Ana}

  \vspace{0.3em}
  JSON esperado:\\
  \texttt{\{"texto":"Ola, Ana!"\}}
\end{frame}

% =============================================================================
\begin{frame}[fragile]{Jeito 1: colocar na URL na mão}
  Vocês podem montar a URL juntando texto:

\begin{lstlisting}[basicstyle=\ttfamily\footnotesize]
var nome = "Ana";

axios.get("/api/saudacao?nome=" + nome)
  .then(function (resposta) {
    console.log(resposta.data.texto);
  });
\end{lstlisting}

  \vspace{0.5em}
  Funciona.\\
  Mas fica confuso quando tem \textbf{vários} parâmetros.
\end{frame}

% =============================================================================
\begin{frame}[fragile]{Jeito 2: objeto params do Axios}
  O Axios tem um jeito mais organizado:

\begin{lstlisting}[basicstyle=\ttfamily\footnotesize]
axios.get("/api/saudacao", {
  params: {
    nome: "Ana"
  }
})
  .then(function (resposta) {
    console.log(resposta.data.texto);
  });
\end{lstlisting}

  \vspace{0.4em}
  O Axios transforma isso em:\\
  \texttt{/api/saudacao?nome=Ana}

  \vspace{0.4em}
  \begin{alertblock}{Usem este jeito}
    Menos erro de digitação.\\
    Mais fácil de ler e de acrescentar parâmetros.
  \end{alertblock}
\end{frame}

% =============================================================================
\begin{frame}[fragile]{Vários parâmetros de uma vez}
\begin{lstlisting}[basicstyle=\ttfamily\footnotesize]
axios.get("/api/soma", {
  params: {
    a: 2,
    b: 3
  }
})
  .then(function (resposta) {
    console.log(resposta.data);
  });
\end{lstlisting}

  \vspace{0.4em}
  Vira na prática:

  \begin{center}
    \texttt{/api/soma?a=2\&b=3}
  \end{center}

  \vspace{0.5em}
  No Java:
\begin{lstlisting}[language=Java,basicstyle=\ttfamily\footnotesize]
@RequestParam int a, @RequestParam int b
\end{lstlisting}
\end{frame}

% =============================================================================
\begin{frame}[fragile]{Parâmetro vindo de um input}
  Na página:

\begin{lstlisting}[basicstyle=\ttfamily\footnotesize]
<input id="campo-nome" type="text">
<button id="btn-saudacao">Saudar</button>
<p id="saida"></p>
\end{lstlisting}

  No \texttt{app.js}:

\begin{lstlisting}[basicstyle=\ttfamily\footnotesize]
var botao = document.getElementById("btn-saudacao");
var campo = document.getElementById("campo-nome");
var saida = document.getElementById("saida");

botao.addEventListener("click", function () {
  axios.get("/api/saudacao", {
    params: { nome: campo.value }
  }).then(function (resposta) {
    saida.textContent = resposta.data.texto;
  });
});
\end{lstlisting}
\end{frame}

% =============================================================================
\begin{frame}{Lendo o fluxo com calma}
  \begin{enumerate}
    \item A pessoa digita \texttt{Ana} no input
    \item Clica no botão
    \item Axios chama \texttt{/api/saudacao?nome=Ana}
    \item O Spring lê \texttt{nome} com \texttt{@RequestParam}
    \item Devolve JSON: \texttt{\{"texto":"Ola, Ana!"\}}
    \item O JS pega \texttt{resposta.data.texto} e mostra na tela
  \end{enumerate}

  \vspace{0.6em}
  \begin{block}{Frase para lembrar}
    \texttt{params} no Axios $\leftrightarrow$ \texttt{@RequestParam} no Spring.
  \end{block}
\end{frame}

% =============================================================================
\begin{frame}{GET com parâmetros $\times$ só GET}
  \begin{center}
  \begin{tabular}{ll}
    \toprule
    Situação & Exemplo \\
    \midrule
    Listar tudo & \texttt{GET /api/produtos} \\
    Buscar com filtro & \texttt{GET /api/produtos?nome=Mouse} \\
    Saudação personalizada & \texttt{GET /api/saudacao?nome=Ana} \\
    \bottomrule
  \end{tabular}
  \end{center}

  \vspace{0.7em}
  Parâmetro na URL = informação \textbf{curta} para filtrar\\
  ou personalizar a consulta.

  \vspace{0.4em}
  Mais à frente veremos o \texttt{POST},\\
  usado para \textbf{enviar dados maiores} e criar registros.
\end{frame}

% =============================================================================
\begin{frame}{Fluxo completo}
  \begin{enumerate}
    \item Usuário abre \texttt{http://localhost:8080/}
    \item Spring entrega \texttt{index.html}, CSS, Axios e JS
    \item O JS executa \texttt{axios.get("/api/produtos")}
    \item Spring roda o método do Controller
    \item Volta JSON com a lista
    \item JS cria elementos \texttt{<li>} e mostra na página
  \end{enumerate}

  \vspace{0.7em}
  \begin{alertblock}{Aula 02 em uma frase}
    Página na pasta \texttt{static/} + dados via \texttt{/api/...}.
  \end{alertblock}
\end{frame}

% =============================================================================
\begin{frame}{Ferramenta de detetive: aba Network}
  No navegador: \textbf{F12} ou botão direito $\rightarrow$ Inspecionar.\\
  Aba \textbf{Network} / Rede.

  \vspace{0.5em}
  Vocês conseguem ver:
  \begin{itemize}
    \item Se \texttt{index.html} carregou
    \item Se \texttt{/api/produtos} foi chamado
    \item Se os parâmetros aparecem na URL
    \item Qual JSON voltou
    \item Se deu erro 404 ou 500
  \end{itemize}

  \vspace{0.6em}
  \begin{block}{Dica}
    Antes de pedir ajuda, olhem a aba Network.\\
    Muitas respostas estão ali.
  \end{block}
\end{frame}

% =============================================================================
\begin{frame}{Erros comuns}
  \begin{itemize}
    \item \textbf{axios is not defined} --- esqueceu o
      \texttt{script} do Axios, ou colocou depois do
      \texttt{app.js}
    \item \textbf{Página em branco} --- veja o Console no F12
    \item \textbf{404} --- caminho do \texttt{@GetMapping}
      diferente do \texttt{axios.get}
    \item \textbf{Parâmetro null no Java} --- nome em
      \texttt{params} diferente do \texttt{@RequestParam}
    \item \textbf{JSON sem os campos} --- faltou o \texttt{get}
      na classe
    \item \textbf{Mudança não apareceu} --- reinicie o Spring
  \end{itemize}
\end{frame}

% #############################################################################
\section{Exercícios --- Parte 1}

% =============================================================================
\begin{frame}{Hora de praticar}
  Vamos montar um \textbf{Painel da Turma} em 4 passos.\\
  Cada exercício \textbf{completa} o anterior.

  \vspace{0.5em}
  Usem um projetinho novo no \texttt{start.spring.io}
  ou o \texttt{ola-spring}.

  \vspace{0.5em}
  \begin{block}{O que o painel terá no final}
    \begin{itemize}
      \item Uma mensagem vinda da API
      \item Uma lista de alunos vinda da API
    \end{itemize}
  \end{block}

  \vspace{0.4em}
  Salve $\rightarrow$ reinicie o Spring $\rightarrow$ teste.\\
  Use o navegador e o F12.
\end{frame}

% =============================================================================
\begin{frame}{Visão geral dos 4 passos}
  \begin{enumerate}
    \item Criar a \textbf{página} com os espaços prontos
          e o script do Axios
    \item Criar a classe \texttt{Mensagem} e o endpoint
          \texttt{/api/mensagem}
    \item Ligar o \textbf{botão} ao \texttt{axios.get} da mensagem
    \item Criar a classe \texttt{Aluno}, o endpoint
          \texttt{/api/alunos} e mostrar a lista
  \end{enumerate}

  \vspace{0.7em}
  No final, a mesma página usa \textbf{dois} endpoints.\\
  É o mesmo raciocínio do StockSales.
\end{frame}

% =============================================================================
\begin{frame}[fragile]{Exercício 1 --- A página do painel}
  Criem:

  \vspace{0.3em}
  \texttt{src/main/resources/static/index.html}\\
  \texttt{src/main/resources/static/js/app.js} \\
  {\small\color{black!60}o \texttt{app.js} pode começar vazio}

  \vspace{0.3em}
  Na página, já deixem estes elementos:
\begin{lstlisting}[basicstyle=\ttfamily\footnotesize]
<h1>Painel da Turma</h1>
<button id="btn-mensagem">Buscar mensagem</button>
<p id="saida">Clique no botao...</p>
<h2>Alunos</h2>
<ul id="lista-alunos"></ul>
<script src="https://cdn.jsdelivr.net/npm/axios/dist/axios.min.js"></script>
<script src="/js/app.js"></script>
\end{lstlisting}

  Teste: \texttt{http://localhost:8080/}\\
  A página abre --- ainda sem dados da API.
\end{frame}

% =============================================================================
\begin{frame}[fragile]{Exercício 2 --- Classe Mensagem e API}
  Criem a classe \texttt{Mensagem} com o atributo \texttt{texto},
  construtor e \texttt{getTexto()}.

  \vspace{0.3em}
  No Controller:

\begin{lstlisting}[language=Java,basicstyle=\ttfamily\footnotesize]
@GetMapping("/api/mensagem")
public Mensagem mensagem() {
  return new Mensagem("Ola da API Spring!");
}
\end{lstlisting}

  \vspace{0.4em}
  Teste \textbf{ainda no navegador}, sem JS:\\
  \texttt{http://localhost:8080/api/mensagem}

  \vspace{0.3em}
  Esperado algo como:\\
  \texttt{\{"texto":"Ola da API Spring!"\}}
\end{frame}

% =============================================================================
\begin{frame}[fragile]{Exercício 3 --- Botão busca a mensagem}
  Continuando a \textbf{mesma} página do Exercício 1.

  \vspace{0.3em}
  No \texttt{app.js}:
\begin{lstlisting}[basicstyle=\ttfamily\footnotesize]
var botao = document.getElementById("btn-mensagem");
var saida = document.getElementById("saida");

botao.addEventListener("click", function () {
  axios.get("/api/mensagem")
    .then(function (resposta) {
      saida.textContent = resposta.data.texto;
    });
});
\end{lstlisting}

  \vspace{0.3em}
  Clicou $\rightarrow$ a frase da API aparece no parágrafo.
\end{frame}

% =============================================================================
\begin{frame}[fragile]{Exercício 4 --- Lista de alunos}
  Ainda no \textbf{mesmo} projeto e na \textbf{mesma} página.

  \vspace{0.3em}
  1) Classe \texttt{Aluno} com \texttt{nome} e \texttt{turma}
     + construtor + gets\\
  2) Endpoint:

\begin{lstlisting}[language=Java,basicstyle=\ttfamily\footnotesize]
@GetMapping("/api/alunos")
public List<Aluno> alunos() {
  List<Aluno> lista = new ArrayList<>();
  lista.add(new Aluno("Ana", "Entra21"));
  lista.add(new Aluno("Bruno", "Entra21"));
  lista.add(new Aluno("Carla", "Entra21"));
  return lista;
}
\end{lstlisting}

  3) No \texttt{app.js}, ao carregar a página:
     \texttt{axios.get("/api/alunos")} e preencher
     \texttt{\#lista-alunos} com \texttt{resposta.data}.
\end{frame}

% =============================================================================
\begin{frame}{Painel pronto --- o que vocês montaram}
  Uma única tela, dois pedidos à API:

  \vspace{0.5em}
  \begin{enumerate}
    \item Clique no botão $\rightarrow$ \texttt{/api/mensagem}
    \item Ao abrir a página $\rightarrow$ \texttt{/api/alunos}
  \end{enumerate}

  \vspace{0.7em}
  \begin{alertblock}{Isso é o ensaio do StockSales}
    Depois, no projeto do curso, a ideia é a mesma:\\
    página HTML + Axios + lista vinda do Spring.
  \end{alertblock}
\end{frame}

% =============================================================================
\begin{frame}{Checklist se travar}
  \begin{itemize}
    \item [ ] A pasta é \texttt{src/main/resources/static/}?
    \item [ ] O script do Axios vem \textbf{antes} do \texttt{app.js}?
    \item [ ] O endpoint começa com \texttt{/api/}?
    \item [ ] O Controller tem \texttt{@RestController}?
    \item [ ] A classe de dados tem os métodos \texttt{get}?
    \item [ ] Os dados estão em \texttt{resposta.data}?
    \item [ ] Reiniciei o Spring depois de salvar?
    \item [ ] Olhei o Console e a aba Network no F12?
  \end{itemize}

  \vspace{0.6em}
  Dúvida? Chamem o professor --- mostrem a tela do F12.
\end{frame}

% =============================================================================
\begin{frame}{Intervalo}
  \begin{center}
    \Large Até já --- retomamos às \textbf{20h20}\\[0.8em]
    \normalsize
    Segunda parte: o projeto \textbf{StockSales}\\
    O sistema que vamos evoluir até o fim do curso
  \end{center}
\end{frame}

% #############################################################################
\section{Parte 2 --- O projeto StockSales}

% =============================================================================
\begin{frame}{Bem-vindos à Parte 2}
  A partir de agora, quase toda aula terá:

  \vspace{0.6em}
  \begin{enumerate}
    \item Um pedaço teórico / exercício curto
    \item Evolução do \textbf{mesmo projeto}
  \end{enumerate}

  \vspace{0.8em}
  Objetivo: sentir como é trabalhar em um sistema\\
  de verdade, desde o começo até login, vendas e banco.

  \vspace{0.6em}
  \begin{alertblock}{Nome do projeto}
    \textbf{StockSales} --- vendas com controle de estoque.
  \end{alertblock}
\end{frame}

% =============================================================================
\begin{frame}{Que problema o StockSales resolve}
  Imaginem uma loja simples:

  \vspace{0.5em}
  \begin{itemize}
    \item Tem \textbf{produtos} com preço e quantidade em estoque
    \item Tem \textbf{clientes} cadastrados
    \item Registra \textbf{vendas}
    \item Quando vende, o \textbf{estoque precisa baixar}
  \end{itemize}

  \vspace{0.7em}
  Esse tipo de sistema aparece bastante em empresas:\\
  loja, estoque, atendimento, sistemas internos.
\end{frame}

% =============================================================================
\begin{frame}{As três ideias centrais}
  \begin{columns}[T]
    \begin{column}{0.32\textwidth}
      \begin{block}{Produto}
        Nome\\
        Preço\\
        Estoque
      \end{block}
    \end{column}
    \begin{column}{0.32\textwidth}
      \begin{block}{Cliente}
        Nome\\
        E-mail\\
        Telefone
      \end{block}
    \end{column}
    \begin{column}{0.32\textwidth}
      \begin{block}{Venda}
        Cliente\\
        Itens\\
        Total\\
        Data
      \end{block}
    \end{column}
  \end{columns}

  \vspace{0.9em}
  A parte mais interessante:\\
  \textbf{vender} $\rightarrow$ \textbf{estoque diminui}.
\end{frame}

% =============================================================================
\begin{frame}{Regra de ouro do estoque}
  Ao confirmar uma venda:

  \vspace{0.5em}
  \begin{enumerate}
    \item Verificar se todo item tem estoque suficiente
    \item Se \textbf{faltar} estoque em qualquer item:\\
          \textbf{recusar a venda inteira}
    \item Se estiver ok: gravar a venda e diminuir o estoque
  \end{enumerate}

  \vspace{0.7em}
  \begin{alertblock}{Importante}
    Não vendemos ``metade''.\\
    Ou a venda fecha toda, ou não fecha.
  \end{alertblock}

  Essa regra será implementada em aulas futuras.
\end{frame}

% =============================================================================
\begin{frame}{Telas que o sistema terá}
  \begin{enumerate}
    \item \textbf{Home} --- resumo do sistema
    \item \textbf{Produtos} --- listar e depois cadastrar/editar
    \item \textbf{Clientes} --- cadastro de clientes
    \item \textbf{Nova venda} --- escolher cliente + produtos
    \item \textbf{Histórico de vendas}
    \item \textbf{Login}
    \item \textbf{Relatórios} --- totais, estoque baixo
  \end{enumerate}

  \vspace{0.5em}
  \begin{block}{Hoje}
    Só a base: projeto no ar + home + \textbf{listar produtos}.
  \end{block}
\end{frame}

% =============================================================================
\begin{frame}{Mapa das próximas aulas}
  \begin{center}
  \begin{tabular}{cl}
    \toprule
    Aula & Foco no StockSales \\
    \midrule
    2 & Projeto + página + listar produtos \\
    3 & CRUD produtos e clientes \\
    4 & Vendas + baixa de estoque \\
    5 & Organização do código e validação \\
    6 & Banco de dados \\
    7 & Autenticação / login \\
    8 & Perfis: admin e vendedor \\
    9 & Relatórios, busca, estoque baixo \\
    10 & Fechamento + desafio final \\
    \bottomrule
  \end{tabular}
  \end{center}
\end{frame}

% =============================================================================
\begin{frame}{Arquitetura que vamos usar}
  \begin{center}
  \begin{tikzpicture}[
      box/.style={draw, rounded corners, align=center, minimum width=2.8cm,
                  minimum height=1.05cm, font=\small\bfseries},
      seta/.style={-{Latex}, thick}
    ]
    \node[box, fill=blue!15] (front) {HTML + CSS + JS\\{\normalfont\tiny static/}};
    \node[box, fill=green!15, right=1.4cm of front] (api) {API JSON\\{\normalfont\tiny /api/...}};
    \node[box, fill=yellow!25, right=1.4cm of api] (dados) {Dados\\{\normalfont\tiny memória depois banco}};

    \draw[seta] (front) -- node[above, font=\tiny]{Axios} (api);
    \draw[seta] (api) -- (dados);
    \draw[seta] (dados.south) -- ++(0,-0.55) -| node[below, font=\tiny, pos=0.25]{JSON} (front.south);
  \end{tikzpicture}
  \end{center}

  \vspace{0.8em}
  Página na \texttt{static/}. Dados no Controller. JSON no meio.
\end{frame}

% =============================================================================
\begin{frame}{API inicial de hoje}
  Hoje precisamos de \textbf{um} endpoint:

  \vspace{0.6em}
  \begin{center}
    \Large \texttt{GET /api/produtos}
  \end{center}

  \vspace{0.8em}
  Resposta: lista JSON de produtos\\
  com id, nome, preço e estoque.

  \vspace{0.6em}
  Nas próximas aulas entram criação, edição, exclusão,
  clientes, vendas e login.
\end{frame}

% =============================================================================
\begin{frame}{O que não vamos fazer hoje}
  Para não misturar demais:

  \vspace{0.5em}
  \begin{itemize}
    \item Cadastro e edição de produtos --- Aula 3
    \item Clientes e vendas
    \item Banco de dados
    \item Login / senha
    \item Layout muito elaborado
  \end{itemize}

  \vspace{0.6em}
  \begin{block}{Meta da noite}
    StockSales sobe.\\
    Home abre.\\
    Produtos aparecem vindo da API.
  \end{block}
\end{frame}

% #############################################################################
\section{Mão na massa --- StockSales}

% =============================================================================
\begin{frame}{Passo 1 --- Criar o projeto}
  Em \texttt{https://start.spring.io}:

  \vspace{0.4em}
  \begin{itemize}
    \item \textbf{Project:} Maven
    \item \textbf{Language:} Java
    \item \textbf{Spring Boot:} versão padrão do site
    \item \textbf{Group:} \texttt{com.entra21}
    \item \textbf{Artifact:} \texttt{stocksales}
    \item \textbf{Package:} \texttt{com.entra21.stocksales}
    \item \textbf{Packaging:} Jar
    \item \textbf{Java:} 21
    \item \textbf{Dependency:} Spring Web
  \end{itemize}

  Gerar, baixar, extrair, abrir no VS Code.
\end{frame}

% =============================================================================
\begin{frame}{Passo 2 --- Pasta do frontend}
  Criem:

  \vspace{0.5em}
  \begin{itemize}
    \item \texttt{src/main/resources/static/index.html}
    \item \texttt{src/main/resources/static/css/estilo.css}
    \item \texttt{src/main/resources/static/js/app.js}
  \end{itemize}

  \vspace{0.5em}
  Na \texttt{index.html}, não esqueçam:
  \begin{itemize}
    \item Título \textbf{StockSales}
    \item Uma lista para os produtos
    \item Script do \textbf{Axios} antes do \texttt{app.js}
  \end{itemize}
\end{frame}

% =============================================================================
\begin{frame}[fragile]{Passo 3 --- Classe Produto}
  Pacote: \texttt{com.entra21.stocksales}

\begin{lstlisting}[language=Java,basicstyle=\ttfamily\footnotesize]
public class Produto {
  private Long id;
  private String nome;
  private double preco;
  private int estoque;

  public Produto(Long id, String nome,
                 double preco, int estoque) {
    this.id = id;
    this.nome = nome;
    this.preco = preco;
    this.estoque = estoque;
  }

  public Long getId() { return id; }
  public String getNome() { return nome; }
  public double getPreco() { return preco; }
  public int getEstoque() { return estoque; }
}
\end{lstlisting}
\end{frame}

% =============================================================================
\begin{frame}[fragile]{Passo 4 --- ProdutoController}
\begin{lstlisting}[language=Java,basicstyle=\ttfamily\footnotesize]
@RestController
public class ProdutoController {

  @GetMapping("/api/produtos")
  public List<Produto> listar() {
    List<Produto> produtos = new ArrayList<>();
    produtos.add(new Produto(1L, "Teclado Mecanico", 250.0, 12));
    produtos.add(new Produto(2L, "Mouse Gamer", 120.0, 30));
    produtos.add(new Produto(3L, "Monitor 24", 900.0, 8));
    produtos.add(new Produto(4L, "Headset", 199.9, 15));
    return produtos;
  }
}
\end{lstlisting}
\end{frame}

% =============================================================================
\begin{frame}{Passo 5 --- JS consome /api/produtos}
  No \texttt{app.js}:
  \begin{enumerate}
    \item Selecionar a lista da página
    \item \texttt{axios.get("/api/produtos")}
    \item Usar \texttt{resposta.data}
    \item Para cada produto, criar um item na tela
          com nome, preço e estoque
  \end{enumerate}

  \vspace{0.6em}
  Testem também a API direto:\\
  \texttt{http://localhost:8080/api/produtos}
\end{frame}

% =============================================================================
\begin{frame}{Como vamos trabalhar daqui pra frente}
  \begin{itemize}
    \item Mesmo projeto toda segunda parte da aula
    \item Cada aula acrescenta uma fatia nova
    \item Evitem apagar o que já funciona --- evoluam
    \item Salvemos com frequência
  \end{itemize}

  \vspace{0.7em}
  \begin{block}{Ideia}
    Vocês não estão fazendo exercício solto.\\
    Estão construindo um produto, aula após aula.
  \end{block}
\end{frame}

% =============================================================================
\begin{frame}{Roteiro sugerido até as 22h}
  \begin{enumerate}
    \item Criar o projeto no start.spring.io
    \item Subir e testar que o servidor responde
    \item Criar \texttt{index.html} + CSS mínimo
    \item Criar \texttt{Produto} + \texttt{ProdutoController}
    \item Testar o JSON no navegador
    \item Ligar o Axios e listar na home
  \end{enumerate}

  \vspace{0.5em}
  Se travarem no meio: peçam ajuda cedo.
\end{frame}

% =============================================================================
\begin{frame}{Critérios de aceite}
  Vocês fecharam a Aula 02 se:

  \vspace{0.5em}
  \begin{itemize}
    \item [ ] O StockSales sobe sem erro
    \item [ ] \texttt{/} mostra a página do sistema
    \item [ ] \texttt{/api/produtos} devolve JSON
    \item [ ] A home lista os produtos via Axios
    \item [ ] Vocês conseguem explicar:\\
              ``a página vem do static; os dados vêm da API''
  \end{itemize}
\end{frame}

% #############################################################################
\section{Recapitulando}

% =============================================================================
\begin{frame}{O que aprendemos hoje}
  \begin{enumerate}
    \item Estrutura de um projeto Spring Boot
    \item Diferença entre \textbf{página} e \textbf{API}
    \item JSON como formato de dados
    \item Classe Java com \texttt{get} virando JSON
    \item Axios no JavaScript para consumir a API
    \item Como enviar parâmetros com \texttt{params}
    \item Nascimento do projeto \textbf{StockSales}
  \end{enumerate}
\end{frame}

% =============================================================================
\begin{frame}{Frases para levar}
  \begin{itemize}
    \item A pasta \texttt{static/} guarda a vitrine: HTML/CSS/JS.
    \item A API \texttt{/api/...} entrega dados em JSON.
    \item O JavaScript, com Axios, é a ponte entre a tela e o Spring.
    \item \texttt{params} no Axios conversa com \texttt{@RequestParam}.
    \item Sem banco ainda: lista em memória está ok.
    \item O StockSales vai crescer até o fim do curso.
  \end{itemize}
\end{frame}

% =============================================================================
\begin{frame}{Na próxima aula}
  \begin{block}{Aula 03}
    \begin{itemize}
      \item CRUD completo de \textbf{produtos}
      \item CRUD de \textbf{clientes}
      \item Novas telas e formulários
      \item Ainda em memória, sem banco
    \end{itemize}
  \end{block}

  \vspace{0.8em}
  Tragam o StockSales funcionando, com listagem de produtos.\\
  Qualquer dúvida de hoje, anotem --- começamos revisando.
\end{frame}

% =============================================================================
\begin{frame}
  \begin{center}
    \Huge Obrigado!\\[0.8em]
    \Large Agora é mão na massa no StockSales.\\[1.2em]
    \normalsize Mauricio Duailibi Neto\\
    Turma Java Entra21
  \end{center}
\end{frame}

\end{document}
